\documentclass[UTF8,a4paper,zihao=-4]{ctexart}
\usepackage[margin=23mm]{geometry}
\usepackage{amsmath,amssymb,bm,booktabs,longtable,graphicx}
\usepackage{xcolor,hyperref,fancyhdr}
\hypersetup{colorlinks=true,linkcolor=blue!50!black,urlcolor=blue!50!black,hypertexnames=false}
\pagestyle{fancy}\fancyhf{}\fancyhead[L]{Shift--Rotation 面形重建}
\fancyhead[R]{公式推导与 MATLAB 实现}\fancyfoot[C]{\thepage}
\setlength{\headheight}{15pt}
\newcommand{\ZZ}{\mathcal Z}
\newcommand{\RMS}{\operatorname{RMS}}
\newcommand{\diag}{\operatorname{diag}}
\title{平移--旋转绝对面形测量\\公式推导、可辨识性与函数复现}
\author{依据 Su 等，Optics Letters 37, 3198--3200 (2012)}
\date{2026 年 10 月 1 日}
\begin{document}
\maketitle
\begin{abstract}
本文逐步复现原论文式（2）--（13）：四次固定测量坐标系中的相对面形相减，消去参考面，建立 Zernike 系数与姿态误差组成的线性系统，并用分块 QR 和小矩阵 SVD 求解。提供适用于 $875\times875$、高度单位为米的四个 XYZ 文本输入的 MATLAB R2022a 实现。平移为沿 $+x,+y$ 各 $0.03$ m，物理旋转为顺时针 $1^\circ$，像素间距为 $0.000362$ m。结果是规定低阶基准后的、有限 Zernike 带宽内的绝对面形；power 项不能由本方法标定。模拟验证与未来实验验证明确区分。
\end{abstract}
\setcounter{tocdepth}{1}
\begingroup\small\tableofcontents\endgroup
\newpage
\section{原文核对与模型前提}
原文为 Dongqi Su, Erlong Miao, Yongxin Sui, Huaijiang Yang,
\emph{Absolute surface figure testing by shift-rotation method using Zernike polynomials},
Optics Letters, 37(15), 3198--3200 (2012), DOI: 10.1364/OL.37.003198。
已核对提供的三页 PDF，包括式（1）--（15）和图 1--4。

四次测得的 $W_1,\ldots,W_4$ 不是单独 wafer 的面形，而是参考面、变换后的待测面、测量姿态项之和。分离参考面的前提是：参考面保持稳定，数据始终在同一探测器坐标系中，wafer 移动前后自身形状不变。重力、支撑变形、温漂及不同参考面漂移不会被此模型自动消除。

原文有三处实现时需明确的细节：
\begin{enumerate}
\item 式（2）排印为 $t_{x1}Z_1$；其后正文称倾斜为 $Z_2,Z_3$，式（3）--（9）也使用 $Z_2$。因此本文将该项按 $t_{x1}Z_2$ 修正，并记录此修正。
\item 文中的 $\Delta Z$ 必须解释为有限变换差值；没有除以平移距离，不是微分导数。以实际平移量和角度直接计算，避免小位移近似。
\item 文中指出 power 无法标定。原始 $K+9$ 列并不保证满秩，不能直接将形式解当作唯一物理解。本文显式处理低阶基准，并检验剩余数值秩。
\end{enumerate}

\section{物理坐标、像素坐标与符号}
设文本每行是 $(i_x,i_y,z)$，高度 $z$ 单位 m。默认像素索引为 $1,\ldots,875$，存储矩阵为 $W(i_y,i_x)$；零起点通过参数指定。设中心 $(c_x,c_y)=(438,438)$，半径 $r_p=437$ 像素，则
\begin{equation}
X=(i_x-c_x)h,\quad Y=\sigma_y(i_y-c_y)h,\quad
x=\frac{i_x-c_x}{r_p},\quad y=\frac{\sigma_y(i_y-c_y)}{r_p},
\end{equation}
其中 $h=0.000362$ m/像素，$\sigma_y=1$ 表示像素 $y$ 沿物理 $+Y$，若图像行向下而物理 $+Y$ 向上，应取 $\sigma_y=-1$。默认 $r_ph=0.158194$ m，模型圆直径为 $0.316388$ m。此口径和中心仅由网格假设得到，必须在真实数据到来后按实际有效口径修改，不能据此断言 wafer 的真实直径。

待测面物理位姿从自身坐标 $\bm u$ 变为探测器坐标 $\bm q$：
\begin{equation}
\bm q=Q(\alpha)\bm u+\bm s,\qquad
\bm u=Q(-\alpha)(\bm q-\bm s),\qquad
Q(\beta)=\begin{bmatrix}\cos\beta&-\sin\beta\\\sin\beta&\cos\beta\end{bmatrix}.
\end{equation}
因此，原文 $Z(x+\Delta x,y)$ 中的 $\Delta x$ 是\textbf{函数自变量位移}，与物理平台沿 $+x$ 的移动相反。对用户给定动作，以从测量图这一侧观察的笛卡尔 $+Y$ 向上约定：
\begin{align}
s_x=s_y&=+0.03\ \mathrm m,& \alpha&=-1^\circ\quad(\text{顺时针}),\\
\Delta x_p=\Delta y_p&=-0.03/h=-82.8729281768\ \text{像素},& \beta&=-\alpha=+1^\circ,\\
\delta_x=\delta_y&=\frac{-0.03}{r_ph}\approx-0.189640568.&&
\end{align}
\begin{equation}T_r(x,y)=(x\cos\beta-y\sin\beta,\ x\sin\beta+y\cos\beta).\end{equation}
这不是整像素平移。函数直接在变换坐标上计算多项式，不舍入 $82.8729$，也不对测量面形插值或平移后相减，否则参考面的同像素相消条件会被破坏。若实际观测方向相反或仪器已作镜像，需修改对应坐标符号。四次位置解释为分别相对于原位置的 $x$ 平移、$y$ 平移、中心旋转；若动作累积，需先改为对应的仿射变换。

\section{Zernike 展开与编号}
原文没有完整说明 81 项的编号规范，因此本实现提供明确可复现的编号，不宣称与厂商 Fringe 或 Noll 输出逐项一致。令 $r=\sqrt{x^2+y^2}$，$\theta=\operatorname{atan2}(y,x)$，径向多项式为
\begin{equation}
R_n^m(r)=\sum_{s=0}^{(n-m)/2}
\frac{(-1)^s(n-s)!}{s!\left(\frac{n+m}{2}-s\right)!\left(\frac{n-m}{2}-s\right)!}r^{n-2s}.
\end{equation}
实基函数在连续单位圆上归一化为单位均方：
\begin{equation}
Z_n^0=\sqrt{n+1}R_n^0,\quad
Z_{n,m}^{c}=\sqrt{2(n+1)}R_n^m\cos(m\theta),\quad
Z_{n,m}^{s}=\sqrt{2(n+1)}R_n^m\sin(m\theta).
\end{equation}
按径向阶 $n$ 递增，同阶按 $m=\operatorname{mod}(n,2),\ldots,n$ 递增；$m>0$ 时先余弦后正弦。前 6 项为
\begin{equation}
\begin{split}
Z_1&=1,\quad Z_2=2x,\quad Z_3=2y,\quad Z_4=\sqrt3(2x^2+2y^2-1),\\
Z_5&=\sqrt6(x^2-y^2),\qquad Z_6=2\sqrt6xy.
\end{split}
\end{equation}
81 项包含 $n\le11$ 的全部 78 项，以及 $n=12$ 的 $m=0$ 和 $m=2$ 余弦、正弦项。任意项的 $(n,|m|,\text{类型})$ 均在函数输出 \texttt{modeTable} 中可查。离散遮挡口径上的基函数不必正交，仍应联合最小二乘，不能直接把点积当作系数。

\section{从原文式（2）--（5）到三个差分式}
\subsection{四次观测}
令 $\bm z(\bm q)=[Z_1(\bm q),\ldots,Z_K(\bm q)]$，
$\bm a=[a_1,\ldots,a_K]^{\mathsf T}$，则待测面
\begin{equation} W_t(\bm q)\simeq\bm z(\bm q)\bm a.\tag{原1}\end{equation}
以下各量都在同一探测器坐标 $\bm q=(x,y)$ 表示：
\begin{align}
W_1&=W_r+\sum_{i=1}^K a_iZ_i(x,y)+t_{x1}Z_2+t_{y1}Z_3+p_1Z_4,\tag{原2修正}\\
W_2&=W_r+\sum_{i=1}^K a_iZ_i(x+\delta_x,y)+t_{x2}Z_2+t_{y2}Z_3+p_2Z_4,\tag{原3}\\
W_3&=W_r+\sum_{i=1}^K a_iZ_i(x,y+\delta_y)+t_{x3}Z_2+t_{y3}Z_3+p_3Z_4,\tag{原4}\\
W_4&=W_r+\sum_{i=1}^K a_iZ_i(r,\theta+\beta)+t_{x4}Z_2+t_{y4}Z_3+p_4Z_4.\tag{原5}
\end{align}
$t_{xj},t_{yj},p_j$ 是以相应 Zernike 基为单位的高度系数（m），不是直接的弧度倾角。式（原5）的 $Z_i(r,\theta+\beta)$ 与 $Z_i(T_r(x,y))$ 完全等价。

\subsection{逐项相减}
对 $W_2-W_1$，同像素 $W_r-W_r=0$，逐项整理得到
\begin{align}
D_xW&=W_2-W_1\notag\\
&=\sum_{i=1}^K a_i\underbrace{[Z_i(x+\delta_x,y)-Z_i(x,y)]}_{\Delta Z_{ix}}
+(t_{x2}-t_{x1})Z_2+(t_{y2}-t_{y1})Z_3+(p_2-p_1)Z_4.\tag{原6}
\end{align}
同理定义
\begin{align}
\Delta Z_{iy}&=Z_i(x,y+\delta_y)-Z_i(x,y),\\
\Delta Z_{ir}&=Z_i(T_r(x,y))-Z_i(x,y),\\
\bm b_j&=(t_{x,j+1}-t_{x1},\ t_{y,j+1}-t_{y1},\ p_{j+1}-p_1)^{\mathsf T},
\quad j=1,2,3.
\end{align}
于是三个方程为
\begin{align}
D_xW&=\sum_i a_i\Delta Z_{ix}+[Z_2\ Z_3\ Z_4]\bm b_1,\tag{原7}\\
D_yW&=\sum_i a_i\Delta Z_{iy}+[Z_2\ Z_3\ Z_4]\bm b_2,\tag{原8}\\
D_rW&=\sum_i a_i\Delta Z_{ir}+[Z_2\ Z_3\ Z_4]\bm b_3.\tag{原9}
\end{align}
例如
\begin{equation}
\Delta Z_{4x}=\sqrt3\,[4\delta_xx+2\delta_x^2],
\end{equation}
其中常数 $\delta_x^2$ 不可在有限平移下任意丢弃。$\partial_xZ_i\,\delta_x$ 只是小平移的一阶近似，本实现不使用此近似。

\section{有效口径筛选与 Gx=R 组装}
令三组可用探测器像素集合为
\begin{equation}
\Omega_j=\{\bm q:\ W_1(\bm q),W_{j+1}(\bm q)\text{均有效},\quad
|\bm q|\le1,\quad|T_j\bm q|\le1\},\quad j=1,2,3.
\end{equation}
每一组可以使用不同有效像素，无需把三组都裁成最小交集。对 $\bm q_{jk}\in\Omega_j$，定义
\begin{equation}
A_j(k,i)=Z_i(T_j\bm q_{jk})-Z_i(\bm q_{jk}),\quad
B_j(k,:)=[Z_2(\bm q_{jk}),Z_3(\bm q_{jk}),Z_4(\bm q_{jk})].
\end{equation}
原文的矩阵可清晰写成块形式
\begin{equation}
\underbrace{\begin{bmatrix}
A_1&B_1&0&0\\ A_2&0&B_2&0\\ A_3&0&0&B_3
\end{bmatrix}}_{G}
\underbrace{\begin{bmatrix}\bm a\\\bm b_1\\\bm b_2\\\bm b_3\end{bmatrix}}_{\bm x}
=
\underbrace{\begin{bmatrix}
\operatorname{vec}_{\Omega_1}(W_2-W_1)\\
\operatorname{vec}_{\Omega_2}(W_3-W_1)\\
\operatorname{vec}_{\Omega_3}(W_4-W_1)
\end{bmatrix}}_{\bm R}.\tag{原10--13}
\end{equation}
若 $N_j=|\Omega_j|$，则 $G$ 为 $(N_1+N_2+N_3)\times(K+9)$。原文给出的 $3MN-M-N$ 对应其示意离散相邻平移索引；对本任务非整数 82.8729 像素平移、圆口径及缺失点，必须以实际掩膜计数，不能硬套该行数。

XYZ 文件不要求行排序；读取后按 $(i_y,i_x)$ 放入矩阵，相减和 $G$ 的行使用同一 MATLAB 线性索引顺序。坐标重复或越界时直接报错。高度 0 是有效值，不可作为无效像素标记。若仪器在孔洞内补 0，应传入独立的逻辑有效掩膜；圆外及 NaN 自动剔除。

\section{不可辨识项、活塞扩展与旋转的作用}
\subsection{为什么不能不加约束直接求逆}
$Z_1=1$ 在任何变换下差分都为零，所以原矩阵至少有一列恒零。实际每次采集还可能含任意高度零点 $c_j$。为容纳这种情况，默认扩展为
\begin{equation}
B_j=[Z_1\ Z_2\ Z_3\ Z_4],\qquad
\bm b_j=[c_{j+1}-c_1,\ t_{x,j+1}-t_{x1},\ t_{y,j+1}-t_{y1},\ p_{j+1}-p_1]^{\mathsf T}.
\end{equation}
这给出 12 个差分姿态系数，是对原文 9 系数模型的\textbf{显式扩展}，不是原文已经写出的项。设置 \texttt{IncludePiston=false} 可使用原文的 9 项模式，但它要求差分常数确实可以忽略或已采用一致的基准处理；不能把任意数据分别去低阶后仍沿用未投影的 $G$。

在默认模型中，对 wafer 加任意
\begin{equation}
f(x,y)=c+u x+v y+w(x^2+y^2)
\end{equation}
产生的三个有限差分均可由各组常数、倾斜姿态项表示，因而数据无法区分。例：
\begin{align}
f(x+\delta,y)-f(x,y)&=u\delta+2w\delta x+w\delta^2,\\
f(T_r\bm q)-f(\bm q)&=[u\ v](Q(\beta)-I)\bm q.
\end{align}
故采用基准约束
\begin{equation}\boxed{a_1=a_2=a_3=a_4=0.}\end{equation}
实现删除这四列，只求 $a_5,\ldots,a_K$。$K=81$ 时默认未知数为 $77+12=89$，原文姿态模式为 $77+9=86$。这种处理承认而非恢复缺失的低阶信息；外部曲率、绝对厚度或基准测量才能补充它们。原文未显式处理常数，若把有限平移带来的常数误当作额外绝对信息，可能得到依赖高度零点的伪解。

\subsection{为何加入旋转能恢复像散}
仅用平移，二次多项式 $x^2-y^2$、$2xy$ 的差分是一次多项式，可被倾斜吸收。旋转下，角向阶 $m$ 的余弦、正弦对子满足
\begin{equation}
\begin{bmatrix}Z_c(\theta+\beta)\\Z_s(\theta+\beta)\end{bmatrix}
=\begin{bmatrix}\cos m\beta&-\sin m\beta\\\sin m\beta&\cos m\beta\end{bmatrix}
\begin{bmatrix}Z_c(\theta)\\Z_s(\theta)\end{bmatrix}.
\end{equation}
其差分矩阵的两个奇异值均为 $2|\sin(m\beta/2)|$。像散 $m=2$ 在 $0^\circ,180^\circ$ 时不可识别；在 $1^\circ$ 时因子约为 $0.034905$，虽可识别但对噪声敏感。径向 $m=0$ 的旋转差分始终为零，高阶径向项需由平移约束。实现对完整剩余矩阵检验秩和条件数，不能只凭角度非零宣称所有模式可解。

\section{稳定最小二乘与误差传播}
真实数据含噪声，通常不存在精确的 $G\bm x=\bm R$，求解应为
\begin{equation}
\widehat{\bm x}=\arg\min_{\bm x}\|G\bm x-\bm R\|_2^2.
\end{equation}
实际计算不构造 $(G^{\mathsf T}G)^{-1}$。将每个数据块的增广矩阵 $[G_b\ \bm R_b]$ 与上次保留的三角因子合并做 QR，只保留新三角因子 $T$：
\begin{equation}
\begin{bmatrix}T_{b-1}\\G_b\ \bm R_b\end{bmatrix}=Q_bT_b.
\end{equation}
最后令 $C=T(:,1:p)$、$\bm d=T(:,p+1)$，有
\begin{equation}
\|G\bm x-\bm R\|_2=\|C\bm x-\bm d\|_2.
\end{equation}
增广 QR 保留了残差能量。令 $s_i=\|C(:,i)\|_2$，$S_c=\diag(s_i)$，对列归一化矩阵作 SVD：
\begin{equation}
C S_c^{-1}=U\Sigma V^{\mathsf T},\qquad
\widehat{\bm x}=S_c^{-1}V\Sigma^{-1}U^{\mathsf T}\bm d.
\end{equation}
仅在所有 $\sigma_i>10^{-11}\sigma_1$ 时接受解，否则报秩亏错误。输出条件数为归一化后 $\sigma_1/\sigma_p$。不做静默正则化或伪逆补解。

在 $G$ 固定、已约束且满列秩时，原文式（14）给出
\begin{equation}
G\Delta\bm x=\Delta\bm R,\qquad
\Delta\bm x=G^+\Delta\bm R
\end{equation}
（有不一致噪声时等号按最小二乘理解）。若平移或角度标定也有误差，则一阶关系应为
\begin{equation}
\Delta\bm x\simeq G^+(\Delta\bm R-\Delta G\,\bm x).
\end{equation}
小平移 $A_x\sim\delta_x\partial_x\ZZ$ 而坐标误差 $\Delta A_x\sim\epsilon_x\partial_x\ZZ$，解释了原文式（15）中相对误差约随 $E_{\rm shift}/Q_{\rm shift}$ 变化的趋势。但比例系数依赖表面、模式、角度与噪声，不是适用于所有数据的常数定理。

同一 $W_1$ 被相减三次会引入相关噪声。若四次独立噪声的方差均为 $\sigma^2$，同一有效像素处三个差分的协方差为
\begin{equation}
\Sigma_D=\sigma^2\begin{bmatrix}2&1&1\\1&2&1\\1&1&2\end{bmatrix}.
\end{equation}
本次复现采用原文的普通最小二乘，不输出未经协方差建模的统计置信区间。若未来需要最小方差估计，可使用实际噪声模型对白化后的系统进行广义最小二乘。

\section{面形输出与三种 RMS}
输出原口径内的面形为
\begin{equation}
\widehat W_t(x,y)=\sum_{i=5}^K\widehat a_iZ_i(x,y).
\end{equation}
它是规定基准下的有限阶多项式重建，不包含任意像素高频粗糙度。原始有效口径外输出 NaN，缺失点不以 0 填补。三种 RMS 必须区分：
\begin{enumerate}
\item 面形 RMS：$\sqrt{N^{-1}\sum(\widehat W_t-\overline{\widehat W_t})^2}$，描述面形本身起伏，不是重建误差。有限掩膜下均值未必严格为零。
\item 方程残差 RMS：$\sqrt{(N_1+N_2+N_3)^{-1}\|G\widehat{\bm x}-\bm R\|^2}$，描述测量差分的拟合情况。无真值时只能计算它，不能由此宣称达到某个绝对准确度。
\item 真值误差 RMS：独立已知 $W_{\rm true}$ 时在共同掩膜内计算误差。为公平比较，去除不可识别的活塞、倾斜、power 差异：
\begin{equation}
\bm e=\widehat{\bm w}-\bm w_{\rm true},\quad
\widehat{\bm c}=\arg\min_{\bm c}\|\bm e-B_0\bm c\|^2,\quad
\RMS_e=\sqrt{N^{-1}\|\bm e-B_0\widehat{\bm c}\|^2},
\end{equation}
其中 $B_0=[Z_1\ Z_2\ Z_3\ Z_4]$ 在共同有效像素取值。同时保留不对齐的原始误差 RMS，防止隐藏不一致。
\end{enumerate}
若设置抽样步长 $s>1$，仅对拟合矩阵抽样；面形输出和差分残差仍在全部有效像素计算。默认 $s=1$ 使用全部像素。输出的 $W_1-\widehat W_t$ 仍含第一次测量姿态误差，命名为 \texttt{referenceWithOriginalAlignment}，不得直接作为无姿态误差的纯参考面。

\section{MATLAB 文件、接口和数据约定}
工作目录为用户指定的 \texttt{code} 目录。主要文件如下：
\begin{center}\small
\begin{tabular}{ll}\toprule
文件&用途\\\midrule
\texttt{fun\_ShiftRotation.m}&读取四个文件、分块求解、面形和 RMS 输出\\
\texttt{sr\_zernike.m}&明确编号的实 Zernike 基\\
\texttt{sr\_read\_xyz.m}&XYZ 读取、索引检查、缺失值处理\\
\texttt{run\_experiment.m}&已填用户位移和角度的实验入口\\
\texttt{sr\_make\_simulation.m}&生成四个模拟输入及真值\\
\texttt{demo\_ShiftRotation.m}&875 网格、81 项的完整模拟验证\\
\texttt{test\_ShiftRotation.m}&独立多项式、秩、数据读取等测试\\\bottomrule
\end{tabular}\end{center}
函数基本用法（平移输入为自变量方向）：
\begin{verbatim}
cfg.ShiftPixels = [-0.03 -0.03]/0.000362;
cfg.RotationDeg = 1;
cfg.PixelPitch = 0.000362;
cfg.GridSize = [875 875];
cfg.K = 81;
cfg.IndexBase = 1;
cfg.YAxisSign = 1;
[wafer,out] = fun_ShiftRotation( ...
 'wafer_original.xyz','wafer_xshift.xyz', ...
 'wafer_yshift.xyz','wafer_rotation.xyz',cfg);
\end{verbatim}
不设置位移和旋转参数时函数直接报错，避免隐藏默认实验动作。当前 XYZ 约定无表头、三列数值，可为制表符、空格等可被 \texttt{readmatrix} 识别的分隔符。每个像素一行，允许无序行和缺失坐标；无效高度写 NaN。若后续表头是固定行数，设置 \texttt{HeaderLines}；若增加列、改变坐标单位或导出为仪器专用格式，则只需先适配读取函数，不能直接错位读取。

输入第三列已按用户定义为表面高度 m，因此没有额外除以 2 或乘波长。若真实设备导出的是反射波前 OPD 或 waves，需在读取层按设备定义转换为表面高度，并确认面形正负号与本文模型一致。

\section{模拟验证记录与实验边界}
模拟输入明确标注为 \texttt{sim\_*.xyz}，不冒充实验数据。低阶测试使用独立写出的笛卡尔多项式作为前向真值，包括像散、彗差、三叶和球差；任意固定参考面另含正弦项。81 项测试额外使用相同 Zernike 基合成高阶项，用于检验精确模型下的实现一致性，因此不是对真实模型失配的证明。模拟每次测量加入独立高斯高度噪声，标准差为 $0.1$ nm。

在本机 MATLAB R2022a 中实际执行了测试及完整模拟，随机种子为 20261001，结果如下（高度误差均为 nm）：
\begin{center}\small
\begin{tabular}{lr}\toprule
检查项目 & 实测运行结果\\\midrule
独立多项式、无噪声，65 网格：对齐误差 RMS & $1.4966\times10^{-14}$\\
81 项、无噪声，101 网格：对齐误差 RMS & $6.2681\times10^{-15}$\\
81 项、0.1 nm 噪声，101 网格：对齐误差 RMS & $0.0805975$\\\midrule
完整模拟网格与项数 & $875\times875,\quad K=81$\\
完整模拟有效系统行数与列数 & $1\,655\,053\times89$\\
完整模拟数值秩 & 89\\
列归一化后条件数 & 37.3141\\
完整模拟真值对齐误差 RMS & $0.00505495$\\
完整模拟未对齐误差 RMS & $0.00505495$\\
完整模拟差分方程残差 RMS & $0.14138460$\\
恢复面形自身 RMS & $1.11495505$\\
生成数据及求解耗时（本机本次） & $38.39$ s\\\bottomrule
\end{tabular}\end{center}
残差接近 $\sqrt2\times0.1=0.141421$ nm，与独立噪声相减的量级一致。较多像素对有限数量的参数提供重复约束，使重建真值误差小于单像素输入噪声；这不保证真实仪器具有相同表现。上述近机器精度的无噪声误差只说明实现一致性，不是物理测量准确度。

测试还核对了原文 9 姿态项模式、0 起点且乱序并带表头的文本、不同有效掩膜、低阶基准对齐、分块 QR 与显式反斜杠求解的一致性。$0^\circ$ 和 $180^\circ$ 的退化旋转应被秩检查拒绝。小网格测试保持与 875 网格相同的归一化位移；只有 875 网格模拟严格对应给定的 $0.03$ m 平移和像素间距。

\IfFileExists{simulation/validation.png}{
\begin{figure}[htbp]\centering
\includegraphics[width=\textwidth]{simulation/validation.png}
\caption{完整模拟的真值、恢复面形及去不可辨识低阶后的误差，单位 nm；不是实验数据。}
\end{figure}}{}

原文报告五次标定的逐像素样本标准差图，其均值约为 $0.13$ nm；与 two-sphere 方法的逐像素差异约为 $0.2$ nm RMS。这两个量分别代表重复性与方法间一致性，本次无真实数据，不能声称复现了这两个实验数字。若重复测量 $L$ 次，应先在共同口径对齐相同基准，再计算
\begin{equation}
s(x,y)=\sqrt{\frac1{L-1}\sum_{\ell=1}^L[W^{(\ell)}(x,y)-\overline W(x,y)]^2},
\quad \overline s=\frac1N\sum_{\Omega}s(x,y).
\end{equation}
它不同于单次面形 RMS，也不同于均值的标准误差 $s/\sqrt L$。未来实验验证应使用独立的 wafer 参考结果或校准方法，在统一坐标、口径及相同低阶基准下计算误差 RMS；没有独立真值时仅报告残差和重复性。

\section{适用范围与待实验数据确认事项}
本实现适用于与式（2）--（5）一致的圆口径刚性面形测量。已采用用户确认的移动量、角度及像素间距。仍需由真实 XYZ 或仪器设置确认：像素起点、y 轴方向、旋转中心、有效口径、文件表头、缺失值规则、输入究竟是高度还是 OPD，以及观察方向是否与旋转符号一致。特别是 $1^\circ$ 旋转对二次角向分量的信号较弱，应关注输出条件数及独立误差验证。中心偏差、角度误差和非刚性变形不在当前线性未知量中，需通过标定或扩展模型处理。
\end{document}
